\chapter{Revisão de Literatura}
\label{cap:revisao-literatura}

\section{Fundamentos}
\label{sec:fundamentos}

% Tópicos sugeridos:
%   - representação SMILES (sintaxe, validação, canonização)
%   - noções de espectroscopia / propriedades ópticas de interesse
%   - TD-DFT e estados excitados (o que é, por que é usado para prever
%     propriedades fotoativas)
%   - representação de moléculas como grafos/árvores

\section{Estado da arte em geração molecular}
\label{sec:estado-da-arte}

% Tópicos sugeridos:
%   - algoritmos genéticos (GA) aplicados à geração/otimização molecular
%   - abordagens baseadas em grafos (GA-grafo, edição de grafos)
%   - modelos de linguagem sobre SMILES (RNN, transformers)
%   - autoencoders variacionais (VAE) para geração molecular
%   - busca em árvore / MCTS aplicada à construção molecular
%   - métricas de similaridade/diversidade estrutural (Tanimoto,
%     fingerprints, distância de edição de grafos/árvores, Zhang-Shasha)

\section{Síntese e lacunas}
\label{sec:sintese-lacunas}

% Tópicos sugeridos:
%   - quadro comparativo dos trabalhos discutidos
%     (abordagem x controle de geração x interpretabilidade x custo)
%   - lacuna que este trabalho pretende preencher
